\section*{الفصل \ref{ch:b2:microbial-metabolism} --- الاستقلاب الجرثومي والأغشية الحيوية}

\begin{solution}{exo:b2:microbial-metabolism:1}
الزُّراق: ذاتي التغذي ضوئيًا معدنيًا (الضوء، والماء،
و$\mathrm{CO_2}$). و\emph{Nitrosomonas}: ذاتي التغذي كيميائيًا
معدنيًا (أكسدة الأمونيوم، و$\mathrm{CO_2}$). و\emph{E.~coli} على
الغلوكوز: غير ذاتي التغذي كيميائيًا عضويًا. والبكتيريا الأرجوانية
غير الكبريتية على السكسينات في الضوء: غير ذاتية التغذي ضوئيًا
عضويًا. و\emph{Thiobacillus} على الكبريتيد في الظلام: ذاتي التغذي
كيميائيًا معدنيًا.
\end{solution}

\begin{solution}{exo:b2:microbial-metabolism:2}
الأكسجين (ميليمترات السطح)، ثم النترات (تحتها مباشرة، حيث زال
الأكسجين)، ثم أكاسيد المنغنيز والحديد (الانتقال من البني إلى
الرمادي)، ثم الكبريتات (الطبقة الكبريتيدية السوداء، وهي سميكة في
الرسابات البحرية)، ثم $\mathrm{CO_2}$ \omterm{def:b2:microbial-metabolism:anaerobic}{لتوليد الميثان} (الأعمق، أو
تحت السطح مباشرة في ماء عذب فقير بالكبريتات).
\end{solution}

\begin{solution}{exo:b2:microbial-metabolism:3}
$5\,\mathrm{C_6H_{12}O_6} + 24\,\mathrm{NO_3^-} + 24\,\mathrm{H^+}
\to 30\,\mathrm{CO_2} + 12\,\mathrm{N_2} + 42\,\mathrm{H_2O}$: أي 4.8
شاردة نترات لكل غلوكوز (فكل نترات تتقبل 5 إلكترونات، وكل غلوكوز
يعطي 24).
\end{solution}

\begin{solution}{exo:b2:microbial-metabolism:4}
جماعة جرثومية ملتصقة بسطح ومغروسة في حشوة من عديدات السكريد
والبروتينات وDNA تصنعها بنفسها. والمراحل: الالتصاق، ثم المستعمرة
المجهرية وإفراز الحشوة، ثم النضج بالأبراج والقنوات، ثم التشتت.
واللويحة السنية والنخر؛ وأخماج القثاطر والغرسات (والرئتان في التليف
الكيسي)؛ وتوسخ الأنابيب والهياكل وتآكلها --- أو، على وجه النفع،
المرشحات الراشحة في معالجة المياه.
\end{solution}

\begin{solution}{exo:b2:microbial-metabolism:5}
تسقط ثمانية إلكترونات من \qty{-0.42}{V} إلى \qty{-0.24}{V}:
$\Delta G^{\circ\prime} = -8\times 96.5\times 0.18 =
\qty{-139}{kJ/mol}$ لكل ميثان (والقيمة المجدولة \qty{-131}{kJ/mol})،
أي \qty{35}{kJ} لكل $\mathrm{H_2}$: أي 0.7 من ATP لكل هدروجين على
الأكثر، ونحو نصف واحد عمليًا. والطاقة لكل جزيء ركيزة صغيرة إلى حد
يجعل مولّدات الميثان تنمو ببطء (أزمنة جيل من ساعات إلى أيام) وتحوّل
جل ركيزتها إلى غاز لا إلى خلايا.
\end{solution}

\begin{solution}{exo:b2:microbial-metabolism:6}
$\mu = 1.2\,S/(5 + S)$: أي \qty{0.20}{h^{-1}} ($T = \ln 2/\mu =
\qty{3.5}{h}$)، و\qty{0.60}{h^{-1}} (\qty{1.2}{h})، و\qty{1.09}{h^{-1}}
(\qty{0.64}{h}). وعند تسعين في المئة: $S = 9K_s = \qty{45}{mg/L}$.
\end{solution}

\begin{solution}{exo:b2:microbial-metabolism:7}
$S^* = 0.05\times 0.4/(0.8 - 0.4) = \qty{0.05}{g/L}$؛ و$X^* =
0.4\times(5 - 0.05) = \qty{1.98}{g/L}$؛ والإنتاج $DX^* = 0.4\times
1.98 = \qty{0.79}{g/L/h}$؛ ومعدل الغسل $D_c = 0.8\times 5/5.05 =
\qty{0.79}{h^{-1}}$.
\end{solution}

\begin{solution}{exo:b2:microbial-metabolism:8}
المقبوض: $0.25\times 275 = \qty{69}{kJ}$، أي 1.4 من ATP لكل
أمونيوم؛ ومنه $12/1.4 = 8.7$، أي نحو تسع شوارد أمونيوم لكل كربون
--- وأضعاف ذلك متى حُسبت كلفة صنع NADH بدفع الإلكترونات صعودًا من
النتريت. أما غير ذاتي التغذي فيحصل على كربونه مرجعًا منظمًا سلفًا،
وعلى ثلاثين ATP لكل غلوكوز؛ في حين على المُنَترِت أن يشتري طاقته
وقدرته الإرجاعية معًا من واهب فقير وأن يبني كل كربون من
$\mathrm{CO_2}$، فيحول كسرًا ضئيلًا من الركيزة التي يعالجها إلى
خلايا.
\end{solution}

\begin{solution}{exo:b2:microbial-metabolism:9}
$L_p = \sqrt{2\times 1.2\times 10^{-9}\times 0.2/0.5} =
\sqrt{9.6\times 10^{-10}} = \qty{31}{\micro m}$. ومضاعفة $c_0$ تضرب
$L_p$ في $\sqrt 2$: \qty{44}{\micro m}؛ وتربيع $q$ ينصّفه:
\qty{15}{\micro m}.
\end{solution}

\begin{solution}{exo:b2:microbial-metabolism:10}
من الأعلى إلى الأسفل: ماء فيه طحالب وزراقم (تركيب ضوئي منتج
للأكسجين)؛ ثم نطاق صدئ من مؤكسدات الحديد حيث يلتقي $\mathrm{Fe^{2+}}$
بالأكسجين؛ ثم حجاب أبيض من مؤكسدات الكبريت عديمة اللون
(\omterm{def:b2:microbial-metabolism:trophic}{متغذيات كيميائية معدنية}) حيث يلتقي الكبريتيد بالأكسجين؛ ثم نطاق
أرجواني من بكتيريا كبريتية أرجوانية (\omterm{def:b2:microbial-metabolism:anoxygenic}{تركيب ضوئي غير منتج للأكسجين}
على الكبريتيد)؛ ثم نطاق أخضر من بكتيريا كبريتية خضراء (الشيء نفسه،
تحتمل كبريتيدًا أكثر وضوءًا أقل)؛ ثم وحل أسود من مخمّرات ومرجعات
كبريتات. (أ) تحتاج البكتيريا الكبريتية الأرجوانية إلى الضوء من
الأعلى والكبريتيد من الأسفل معًا، فتقع عند السطح البيني حيث يتقاطع
التدرجان. (ب) وتصنع مرجعات الكبريتات $\mathrm{H_2S}$ الذي يرسّب
الحديد كبريتيدًا أسود FeS. (ج) وهي مغلقة من جهة المادة: فالكبريت
والكربون والنتروجين تدور بين صور مؤكسدة ومرجعة دون أن تخرج؛ ومفتوحة
من جهة الطاقة: فالضوء يدخل والحرارة تخرج، ودون ضوء تتوقف كل دورة.
\end{solution}

\begin{solution}{exo:b2:microbial-metabolism:11}
الإنتاج $pn$، والفقد $kA$: ومنه $A^* = pn/k$. و$A_c = 10^{-8}\times
6.02\times 10^{23} = 6.0\times 10^{15}$ جزيئة في اللتر؛ فيكون $n_c =
kA_c/p = 5\times 10^{-4}\times 6.0\times 10^{15}/5 = 6\times 10^{11}$
في اللتر، أي $6\times 10^{8}$ في المليلتر. فالمزرعة الكثيفة
($10^{9}$) فوق العتبة، وماء البحر ($10^{5}$) دونها بأربع رتب من
القدر. وفي \omterm{def:b2:microbial-metabolism:biofilm}{الغشاء الحيوي} تقع الخلايا عند $10^{11}$ في مليلتر من
الحشوة، وتبطئ الحشوة إفلات الإشارة (فتخفض $k$ فعليًا)، فتبلغ
مستعمرة مجهرية من بضعة آلاف خلية في نانولتر النصابَ، بينما لا تبلغه
الخلايا نفسها أبدًا لو تشتتت في لتر.
\end{solution}

\begin{solution}{exo:b2:microbial-metabolism:12}
تثبيت النتروجين (فالنتروجيناز لا يوجد إلا عند بدائيات النوى): لو
توقف لنزف النتروجين المؤتلف في المحيط الحيوي \omterm{def:b2:microbial-metabolism:anaerobic}{بنزع النترتة} وانهار
الإنتاج في غضون عقود. والنترتة \omterm{def:b2:microbial-metabolism:anaerobic}{ونزع النترتة} (عند بدائيات النوى
وحدها): لو توقفتا لتراكم الأمونيوم وزالت النترات، ولما عاد
النتروجين إلى الهواء أبدًا. \omterm{def:b2:microbial-metabolism:anaerobic}{وتوليد الميثان} \omterm{def:b2:microbial-metabolism:anaerobic}{والتنفس اللاهوائي} (عند
بدائيات النوى وحدها): لو توقفا لتراكمت المادة العضوية في الرسابات
اللاأكسجينية والكروش والمستنقعات دون تمعدن ولغادر الكربون الدورة؛
وكذلك كان التركيب الضوئي المنتج للأكسجين ابتكارًا بكتيريًا، ولولاه
لما كان ثمة أكسجين يُتنفس. فحقيقيات النوى إضافة متأخرة وضيقة
كيميائيًا إلى كوكب لبدائيات النوى.
\end{solution}

\begin{solution}{pb:b2:microbial-metabolism:1}
\textbf{1.} $-24\times 96.5\times (0.82 + 0.43) = \qty{-2900}{kJ/mol}$.
\textbf{2.} النترات: $-24\times 96.5\times 1.17 = \qty{-2700}{kJ/mol}$؛
والكبريتات: $-24\times 96.5\times 0.21 = \qty{-490}{kJ/mol}$.
\textbf{3.} يعطي الأكسجين أكثر، فتغلب الهوائيات نازعات النترتة على
المادة العضوية ما دام الأكسجين موجودًا، والأكسجين يكبت مختزلات
النترات؛ أما \omterm{def:b2:microbial-metabolism:anaerobic}{إرجاع الكبريتات} فيحتاج إلى انعدام الأكسجين \emph{وإلى}
غياب النترات، فرائحته تعني أن الحوض ناقص التهوية والنترات مستنفدة
--- و$\mathrm{H_2S}$ سام ويأكل الخرسانة.
\textbf{4.} $0.4\times 2900/50 = 23$ من ATP؛ و$0.4\times 2700/50 = 22$؛
و$0.4\times 490/50 = 4$.
\textbf{5.} نحو ست مرات أكثر ($23/4$).
\textbf{6.} جزيئان من ATP لكل غلوكوز يعنيان مردودًا نحو \num{0.1}،
فتبقى \qty{90}{\%} من طاقة الركيزة في النواتج (حموض، وكحول،
و$\mathrm{H_2}$)، وهي التي تحولها مولّدات الميثان إلى ميثان بكسب
ضئيل مثله؛ فطاقة قليلة، وكتلة حيوية قليلة، والكربون يغادر غازًا.
\textbf{7.} $10^4\times 0.3 = \qty{3000}{kg}$ من COD في اليوم.
\textbf{8.} الكتلة الحيوية $0.4\times 3000 = \qty{1200}{kg}$ من COD؛
أما \qty{1800}{kg} الباقية من COD فتُؤكسد وتحتاج إلى \qty{1800}{kg}
من $\mathrm{O_2}$.
\textbf{9.} $1800/2 = \qty{900}{kWh}$ في اليوم.
\textbf{10.} $10^4\times 0.04 = \qty{400}{kg}$ من النتروجين؛
و$4.57\times 400 = \qty{1830}{kg}$ من $\mathrm{O_2}$.
\textbf{11.} في الحالة المستقرة $\mu = 1/\theta \le \mu_{\max}$:
فالغسل يقع دون $\theta = 1/\mu_{\max} = \qty{2}{d}$.
\textbf{12.} $D = \qty{0.1}{d^{-1}}$: ومنه $S^* = 1\times 0.1/(0.5 - 0.1)
= \qty{0.25}{g/m^3}$.
\textbf{13.} $S^* = 0.1/(0.25 - 0.1) = \qty{0.67}{g/m^3}$، ويرتفع حد
الغسل إلى \qty{4}{d}؛ وعلى المشغّل أن يطيل عمر الحمأة (بطرح حمأة
أقل): فعند $\theta = \qty{20}{d}$ يعود $S^* =
0.05/0.2 = \qty{0.25}{g/m^3}$.
\textbf{14.} $2.86\times 400 = \qty{1144}{kg}$ من COD في اليوم.
\textbf{15.} أمونيوم الصبيب $0.25\times 10^4 = \qty{2.5}{kg/d}$، أي
\qty{0.6}{\%} من الحمولة؛ أما الباقي، وهو \qty{397}{kg} من
النتروجين، فيغادر بصيغة $\mathrm{N_2}$: أي \qty{99}{\%}.
\textbf{16.} الموفَّر: \qty{1144}{kg} من $\mathrm{O_2}$. والطلب
الكلي: $(1800 - 1144) + 1830 = \qty{2490}{kg}$ من $\mathrm{O_2}$ في اليوم.
\textbf{17.} $1200\times 0.6\times 0.35 = \qty{252}{m^3}$ من الميثان
في اليوم.
\textbf{18.} $252\times 36 = \qty{9070}{MJ} = \qty{2520}{kWh}$؛
والكهرباء $0.35\times 2520 = \qty{880}{kWh}$ في اليوم.
\textbf{19.} التهوية $2490/2 = \qty{1245}{kWh}$ في اليوم: فالميثان
يغطي منها $880/1245 = \qty{71}{\%}$.
\textbf{20.} $D_e\,c'' = q$، ومنه $c = c_0 + ax + qx^2/2D_e$؛ ومع
$c(L_p) = 0$ و$c'(L_p) = 0$: $a = -qL_p/D_e$ و$c =
(q/2D_e)(L_p - x)^2$، مع $L_p^2 = 2D_ec_0/q$.
\textbf{21.} $L_p = \sqrt{2\times 1.5\times 10^{-9}\times 0.2/0.3} =
\sqrt{2\times 10^{-9}} = \qty{45}{\micro m}$.
\textbf{22.} $45/300 = \qty{15}{\%}$ هوائي. وتحته ترجع نازعات
النترتة النتراتِ المنتشرة نزولًا من الطبقة السطحية المُنَترِتة إلى
$\mathrm{N_2}$، مستعملة المادة العضوية المنتشرة إلى الداخل.
\textbf{23.} داخل غشاء واحد ينشئ تدرج الأكسجين طبقة هوائية (نترتة)
فوق طبقة لاأكسجينية (نزع نترتة)، تفصلهما بضع عشرات من
الميكرومترات؛ أما الحوض المهوّى فأكسجيني في كل مكان، فيحتاج \omterm{def:b2:microbial-metabolism:anaerobic}{نزع
النترتة} إلى حوض لاأكسجيني منفصل.
\textbf{24.} يتفاعل الكلور مع بوليمرات الحشوة وتستهلكه قبل أن يبلغ
الخلايا العميقة؛ والخلايا العميقة جائعة بطيئة النمو أو ساكنة وأقل
حساسية بكثير؛ فلا تموت إلا الطبقة السطحية وينمو الغشاء من جديد من
الأسفل.
\textbf{25.} طلب الأكسجين \qty{2490}{kg} في اليوم؛ والميثان
\qty{252}{m^3} في اليوم؛ ويوفر الميثان \qty{71}{\%} من كهرباء
التهوية (\qty{880}{kWh} من \qty{1245}{kWh}).
\end{solution}
